\chapter{Gestión del proyecto}
\label{chap:gestion}

\section{Ciclo de vida}

El proyecto ha seguido un ciclo de vida incremental, con siete iteraciones en las que cada una
entrega algo utilizable y se valida antes de pasar a la siguiente. No es una elección de estilo: en
un trabajo con tanta incertidumbre sobre el mecanismo y los componentes, cada iteración ha servido
para descartar una hipótesis. Las fases de definición y análisis se concentraron en la primera
iteración; el diseño, la construcción y las pruebas se repiten en todas.

La validación se ha apoyado sistemáticamente en programas de prueba por subsistema
(sección~\ref{sec:hw-validacion}), de modo que ningún fallo de integración obligara a depurar todo
el sistema a la vez.

\section{Planificación}

\subsection{Planificación real}

La planificación se ha reconstruido a partir del historial del repositorio: 36 \emph{commits} entre
el 4 de marzo y el 21 de septiembre de 2026, concentrados en nueve semanas
(figura~\ref{fig:actividad-git}). El historial sólo refleja el trabajo sobre código y diseño
electrónico; los periodos sin \emph{commits} corresponden a montaje mecánico, pruebas sin cambios de
código y periodo lectivo \completar{describir las actividades de abril a junio y de julio a agosto}.

\begin{figure}[tbp]
\centering
\includegraphics[width=\textwidth]{actividad_git.pdf}
\caption{\emph{Commits} por semana en el repositorio del proyecto.}
\label{fig:actividad-git}
\end{figure}

\begin{description}[style=nextline,leftmargin=1.5em]
  \item[F1. Exploración del accionamiento (2--10 de marzo)] Comunicación Modbus con una placa XIAO
        ESP32-S3: escaneo de direcciones, modos par y velocidad, monitor de terminal, parámetros del
        accionamiento y resistencia de frenado. Culmina en un firmware v1.0 con regulador de fuerza
        en modo par.
  \item[F2. Diseño de la placa adaptadora (12--17 de marzo)] Esquema y colocación de la revisión~A.
  \item[F3. Migración a Raspberry Pi Pico 2 y validación de la placa (20--30 de junio)] Programas de
        prueba de entradas y salidas, células de carga, RS-485 y escaneo Modbus. Versión de doble
        núcleo.
  \item[F4. Firmware de posición en modo velocidad (4 de julio)] Movimiento entre extremos con
        integración de la posición y monitor asociado.
  \item[F5. Firmware v2 y supervisión web (2--6 de septiembre)] \emph{Homing} con dos finales de
        carrera, servidor FastAPI, página de supervisión, ciclos automáticos y grabación de ensayos.
  \item[F6. Primera campaña de ensayos (14--15 de septiembre)] Quince ensayos con el ciclo
        A$\rightarrow$B$\rightarrow$A, visor de ensayos con curva teórica y primer análisis.
  \item[F7. Rediseño del ciclo y del modelo, y segunda campaña (21 de septiembre)] Ciclo
        B$\rightarrow$A$\rightarrow$B con extremos por finales de carrera, modelo con la salida del
        cable desplazada, configuración única del mecanismo con simulador, segunda campaña de
        ensayos y análisis definitivo.
\end{description}

\begin{figure}[tbp]
\centering
\resizebox{\textwidth}{!}{%
\begin{ganttchart}[
    x unit=0.37cm, y unit chart=0.52cm, y unit title=0.55cm,
    vgrid={*1{black!8, dotted}}, hgrid,
    title label font=\scriptsize, bar label font=\scriptsize,
    milestone label font=\scriptsize\itshape,
    bar/.append style={fill=blue!45}, milestone/.append style={fill=red!70},
    bar height=0.55, group label font=\scriptsize\bfseries,
    inline=false
  ]{1}{31}
  \gantttitle{Mar}{5} \gantttitle{Abr}{4} \gantttitle{May}{4} \gantttitle{Jun}{5}
  \gantttitle{Jul}{4} \gantttitle{Ago}{5} \gantttitle{Sep}{4} \\
  \ganttbar{F1 Exploración accionamiento}{1}{2} \\
  \ganttmilestone{Firmware v1.0}{2} \\
  \ganttbar{F2 Diseño PCB rev. A}{2}{3} \\
  \ganttbar[bar/.append style={fill=black!15}]{Actividad fuera de git}{4}{15} \\
  \ganttbar{F3 Migración a Pico 2 y pruebas}{16}{18} \\
  \ganttbar{F4 Firmware de posición}{18}{18} \\
  \ganttbar[bar/.append style={fill=black!15}]{Actividad fuera de git}{19}{26} \\
  \ganttbar{F5 Firmware v2 y web}{27}{28} \\
  \ganttmilestone{Primer ciclo grabado}{28} \\
  \ganttbar{F6 Primera campaña}{29}{29} \\
  \ganttbar{F7 Ciclo y modelo nuevos, 2.ª campaña}{30}{30} \\
  \ganttbar{Redacción de la memoria}{29}{31}
\end{ganttchart}%
}
\caption{Diagrama de Gantt reconstruido (semanas desde el 2 de marzo de 2026).}
\label{fig:gantt}
\end{figure}

\section{Presupuesto}

\subsection{Personal}

La dedicación de un TFG se fija en créditos ECTS, cada uno de 25 a 30 horas de trabajo
\cite{rd1125}. La tabla~\ref{tab:horas} reparte una dedicación estimada por fases \verificar{ajustar
las horas y los créditos del TFG en la titulación}. Para valorar el trabajo se toma como supuesto un
coste de 30~€/h para un ingeniero o ingeniera junior y de 45~€/h para la tutoría.

\begin{table}[H]
\centering
\caption{Dedicación estimada y coste de personal.}
\label{tab:horas}
\begin{tabular}{@{}lrrr@{}}
\toprule
Actividad & Horas & €/h & Coste (€) \\
\midrule
Estudio y documentación & 30 & 30 & 900 \\
Exploración del accionamiento (F1) & 35 & 30 & 1.050 \\
Diseño de la placa (F2) & 45 & 30 & 1.350 \\
Pruebas de placa y comunicaciones (F3) & 45 & 30 & 1.350 \\
Firmware (F4--F5) & 70 & 30 & 2.100 \\
Software de supervisión (F5, F7) & 60 & 30 & 1.800 \\
Ensayos y análisis (F6, F7) & 40 & 30 & 1.200 \\
Redacción de la memoria & 50 & 30 & 1.500 \\
\midrule
Subtotal trabajo técnico & 375 & & 11.250 \\
Tutoría & 25 & 45 & 1.125 \\
\midrule
\textbf{Total recursos humanos} & & & \textbf{12.375} \\
\bottomrule
\end{tabular}
\end{table}

\subsection{Material}

\begin{longtable}{@{}L{0.42\textwidth}rrr@{}}
\caption{Material del sistema.}\label{tab:material}\\
\toprule
Elemento & Cant. & Precio unit. (€) & Total (€) \\
\midrule
\endfirsthead
\toprule
Elemento & Cant. & Precio unit. (€) & Total (€) \\
\midrule
\endhead
Raspberry Pi Pico 2 & 1 & \completar{} & \\
Transceptor RS-485 THVD1400 & 1 & \completar{} & \\
Interruptor de lado alto TPS4H000-Q1 & 1 & \completar{} & \\
Amplificador operacional MCP6294 & 1 & \completar{} & \\
Diodos zéner BZX384 & 4 & \completar{} & \\
Resistencias y condensadores SMD 0805 & 25 & \completar{} & \\
Bornas de tornillo y tiras de pines & -- & \completar{} & \\
Fabricación de la PCB & 1 & \completar{} & \\
Load Cell 2 Click (ZSC31014) & 1 & \completar{} & \\
Célula de carga & 1 & \completar{} & \\
Servomotor paso a paso de lazo cerrado con accionamiento & 1 & \completar{} & \\
Fuente de alimentación del accionamiento & 1 & \completar{} & \\
Finales de carrera & 2 & \completar{} & \\
Caracola, barra, cable, estructura y masa de \SI{2}{\kilo\gram} & -- & \completar{} & \\
\midrule
\textbf{Total material} & & & \completar{} \\
\bottomrule
\end{longtable}

\subsection{Equipos y software}

Todo el software empleado es libre o gratuito (KiCad, Arduino IDE con el núcleo Arduino-Pico,
Python, FastAPI, Chart.js, git y \LaTeX), por lo que su coste es nulo. Los equipos se imputan por
amortización lineal,
\begin{equation}
  A = C\,\frac{t_{\text{uso}}}{t_{\text{vida}}},
\end{equation}
con $C$ el coste de adquisición, $t_{\text{uso}}$ el tiempo de uso y $t_{\text{vida}}$ la vida útil.

\begin{table}[H]
\centering
\caption{Amortización de equipos.}
\label{tab:amortizacion}
\begin{tabular}{@{}lrrrr@{}}
\toprule
Equipo & Coste (€) & Vida útil (meses) & Uso (meses) & Amortización (€) \\
\midrule
Ordenador portátil & \completar{} & 48 & 7 & \\
Polímetro & \completar{} & 60 & 7 & \\
Fuente de laboratorio & \completar{} & 60 & 7 & \\
\midrule
\textbf{Total} & & & & \completar{} \\
\bottomrule
\end{tabular}
\end{table}

\subsection{Resumen de costes}

\begin{table}[H]
\centering
\caption{Presupuesto total.}
\label{tab:presupuesto}
\begin{tabular}{@{}lr@{}}
\toprule
Concepto & Importe (€) \\
\midrule
Recursos humanos & 12.375 \\
Material & \completar{} \\
Amortización de equipos & \completar{} \\
\midrule
Subtotal & \completar{} \\
Costes indirectos (\SI{10}{\percent}, supuesto) & \completar{} \\
\midrule
Base imponible & \completar{} \\
IVA (\SI{21}{\percent}) & \completar{} \\
\midrule
\textbf{Total} & \completar{} \\
\bottomrule
\end{tabular}
\end{table}

\section{Impacto}

\subsection{Impacto técnico y económico}

La consecuencia práctica del capítulo~\ref{chap:modelo} es de dimensionado: para un recorrido dado,
un perfil de par constante reduce el par de pico que debe dar el motor. Sobre el recorrido con carga
de los ensayos, el perfil ideal rebaja el pico un \SI{11}{\percent} respecto a un cilindro
equivalente, lo que permitiría un motor y un accionamiento algo menores. A velocidad constante, las
pérdidas por efecto Joule son proporcionales a $\int\tau^2\,\mathrm{d}\varphi$, mínima cuando el par
es constante (desigualdad de Cauchy-Schwarz), aunque en este caso la mejora es de apenas un
\SI{1}{\percent}.

Esas cifras deben leerse junto con el resultado del capítulo~\ref{chap:resultados}: el tambor
instalado no aplana el par, sino que lo empeora respecto a un cilindro. El margen de mejora que
ofrece el rediseño del perfil es, por tanto, bastante mayor que ese \SI{11}{\percent}.

La plataforma desarrollada es además reutilizable: la placa, el firmware y la aplicación web forman
un banco genérico para cualquier eje accionado por Modbus con medida de fuerza, con un coste de
material muy inferior al de un sistema de adquisición comercial.

\subsection{Objetivos de Desarrollo Sostenible}

\begin{table}[H]
\centering
\caption{Relación del trabajo con los Objetivos de Desarrollo Sostenible \cite{agenda2030}.}
\label{tab:ods}
\begin{tabularx}{\textwidth}{@{}lX@{}}
\toprule
ODS & Contribución \\
\midrule
4. Educación de calidad & Banco de ensayo abierto y documentado, útil como práctica de laboratorio de mecanismos, accionamientos y adquisición de datos. \\
7. Energía asequible y no contaminante & Un perfil de par uniforme reduce las pérdidas en el motor. \\
9. Industria, innovación e infraestructura & Integración de tecnologías industriales (Modbus, RS-485) con herramientas libres de bajo coste. \\
12. Producción y consumo responsables & Dimensionado de actuadores más pequeños; módulo microcontrolador en zócalo, reutilizable; datos brutos conservados para no repetir ensayos. \\
\bottomrule
\end{tabularx}
\end{table}

\subsection{Aspectos éticos, legales y de seguridad}

El prototipo es un equipo de laboratorio que no se comercializa. Su evolución hacia una máquina de
uso por terceros tendría que considerar la Directiva 2006/42/CE de máquinas, que el Reglamento (UE)
2023/1230 sustituye a partir de enero de 2027, la Directiva 2014/30/UE de compatibilidad
electromagnética y la Directiva 2011/65/UE sobre sustancias peligrosas (RoHS). La
sección~\ref{sec:arq-seguridad} discute las limitaciones de seguridad del diseño actual frente a
ISO 13849-1 \cite{iso13849}.

En cuanto a la ética profesional, se ha priorizado la trazabilidad: se conservan las medidas brutas,
las figuras se generan con programas reproducibles y las limitaciones se señalan de forma explícita,
incluso cuando el resultado contradice la hipótesis de partida.
